\section{The S13 backbone and its original closed calculation}
\label{sec:s13}
S13 starts with total continuous mass $M$ and ending fraction $p_s$ in
preparation state $s$. The other mass is a symmetric continuing group.
A common constant advice $v$ gives traversal time $\tau=L_{13}/v$.
With constant positive completion rate $\lambda_{13}$,
\begin{equation}
 \dot E=-\lambda_{13}E,\quad \dot C=\lambda_{13}E,
 \qquad E(\tau)=p_sM e^{-\lambda_{13}\tau}.
 \label{eq:s13transfer}
\end{equation}
Target receiving is assumed nonbinding during preparation. Remaining E
consumes $\alpha>1$ receiving work per mass at a fixed work rate $\mu$:
\begin{equation}
 W_{13}=M+(\alpha-1)E(\tau),\quad
 T_{13}=W_{13}/\mu,\quad
 \overline J_{13}=\tau+T_{13}/2
 =\tau+\frac{M}{2\mu}+K_{13}p_s e^{-\lambda_{13}\tau},
 \label{eq:s13objective}
\end{equation}
where $K_{13}=(\alpha-1)M/(2\mu)$. The half-clearance term follows
from homogeneous proportional fluid-batch departures. It is not a general
formula for an arbitrary arrival or departure curve.

The original action result follows from
\begin{equation}
 \overline J_{13,\tau}=1-K_{13}\lambda_{13}p_s e^{-\lambda_{13}\tau},
 \qquad
 \overline J_{13,\tau\tau}=K_{13}\lambda_{13}^2p_s e^{-\lambda_{13}\tau}>0.
 \label{eq:s13derivative}
\end{equation}
If two aggregate-history-matched preparations satisfy
\begin{equation}
 K_{13}\lambda_{13}p_L<e^{\lambda_{13}\tau_{\min}}
 <K_{13}\lambda_{13}p_H<e^{\lambda_{13}\tau_{\max}},
 \label{eq:s13split}
\end{equation}
strict convexity gives fastest advice for $L$ and an interior slower
$\tau_H^*=\lambda_{13}^{-1}\log(K_{13}\lambda_{13}p_H)$ for $H$.
The chain rule gives
\begin{equation}
 \frac{d\overline J_{13}}{dv}
 =-\frac{L_{13}}{v^2}
 \left(1-K_{13}\lambda_{13}p_s e^{-\lambda_{13}L_{13}/v}\right).
 \label{eq:s13speed}
\end{equation}
Identical aggregate observations force a deterministic observer-only
controller to take the same action. The two unique minimizers differ; the
maximum of their nonnegative continuous regret functions is positive on
the compact action interval and has a positive minimum. This is S13's
conditional information/action/analytical-gradient certificate.

The refinements below replace the constant completion, grouped receiving,
batch timing and fixed-service reductions. Therefore
\eqref{eq:s13objective}--\eqref{eq:s13speed} are historical reference
results, not formulas for the current model. Their validity under S13's
assumptions is preserved.

\section{Current geometry, notation, units and admissible domain}
\label{sec:domain}
\subsection{Setting and initialized state}
The road has preparation length $L_p$ and terminal length $b>0$,
$L=L_p+b$. At decision time a uniform cohort occupies $[0,\ell]$,
$0<\ell\leq L_p$, and the terminal is empty. E ends at $L$; A is its
adjacent continuing lane and B the far continuing lane. Initial densities
satisfy
\begin{equation}
 e_0,a_0,b_0\geq0,\quad a_0=b_0,\quad
 r=e_0+a_0,\quad\rho_0=e_0+2a_0,\quad M=\ell\rho_0>0.
 \label{eq:initial}
\end{equation}
There are no postdecision external arrivals. The constant common advice
is $v\in[v_-,v_+]$, $0<v_-<v_+$. It sets ideal upstream/free travel;
queued boundary motion need not equal it. H1's remaining upstream occupied
body is homogeneous. Terminal propagation is represented by boundary counts;
there is no new interior spatial grid or independent approach-speed ODE.

Use
\begin{equation}
 K=wk_j,\quad\kc(v)=\frac K{w+v},\quad
 C(v)=\frac{Kv}{w+v},\quad c(u)=\frac{Ku}{w+u},\quad
 U(f)=\frac{wf}{K-f},\quad U(C(v))=v.
 \label{eq:fd}
\end{equation}
The shared A/E budget is $\budget$; $B$ as a lane index still means the
far lane. The archived recovery reading $p$ is written $\loss$ here.
$K_{13}$ in Section~\ref{sec:s13} is unrelated to $K=wk_j$.

\subsection{Retained sufficient domain}
All geometry and physical parameters are common across advice and observers.
Retain H1/T1's initial/handoff conditions, in particular
\begin{align}
 &w,k_j,h,a_r>0,\qquad\gamma\geq1,\qquad
 e_0,a_0,b_0\leq\kc(v_+),\nonumber\\
 &v_+r<R\leq C(v_-),\qquad
 v_+b_0\leq R_B\leq C(v_-),\qquad M<k_jb. \label{eq:domain}
\end{align}
For immediate work set fixed reference units once,
\begin{equation}
 \begin{gathered}
 v_\star=v_+,\quad C_\star=C(v_\star),\quad
 c_0=1+\delta,\quad\chi=\frac\beta{C_\star v_\star},\\
 \delta,\beta\geq0,\quad
 \eta=c_0-\beta(1+v_+/w)>0,\quad c_{\max}=c_0+\beta.
 \end{gathered}
 \label{eq:workdomain}
\end{equation}
The anchor fixes $\chi$ once; it is not recomputed when advice or a search
box changes. The reference mass $m=1$ vehicle-equivalent is fixed physical
normalization, not packet size. Ordinary combined E/A arrival load fits
$R$; any primary queue originates in additional forced work. These bounds
are sufficient and may be conservative. They are not an empirically
admitted operating range.

\begin{center}
\begin{tabular}{>{\raggedright\arraybackslash}p{47mm}>{\raggedright\arraybackslash}p{99mm}}
\toprule
Symbols & Units and meaning\\
\midrule
$L,L_p,b,\ell,h,z$ & m; lengths and traveled progress\\
$e,a,b_0,k_j,\kc,\rho_0$ & veh-eq/m; densities\\
$v,w,u_j,U,V$ & m/s; speeds or specified boundary speed quantities\\
$M,F_j,P_j,Y_j,N_j,q_j,m$ & veh-eq; traffic mass, counts or reference mass\\
$f,g_j,\lambda_j,R,R_B,K,\budget,\loss,D$ & veh-eq/s in normalized receiving-work units\\
$f/h$ & veh-eq/(m s); actual lateral density source\\
$a_r$ & m/s$^2$; selected physical recovery acceleration\\
$\gamma,\delta,\beta,c_j$ & dimensionless work coefficients\\
$\chi$ & s$^2$/(veh-eq m); fixed mismatch coefficient\\
$\Theta,t,T$ & s; duration, time or mass-rank admission time\\
$W$ & veh-eq; integrated isolated lost service\\
$J$ & veh-eq s; complete local TST\\
\bottomrule
\end{tabular}
\end{center}

\section{Preparation: conservative transfer limited by adjacent service}
\label{sec:h1}
\subsection{The local rule and free-branch reduction}
Own A demand claims receiving before the E offer:
\begin{equation}
 S(v,a)=\min\{C(v),w(k_j-a)\},\quad
 f=\min\left\{ve,\frac{S(v,a)-va}{\gamma}\right\},\quad
 \dot e=-f/h,\quad\dot a=f/h,\quad\dot b_0=0.
 \label{eq:h1general}
\end{equation}
This is H1's opportunity-equivalent closure. Actual transfer over
$dx\,dt$ is $(f/h)dx\,dt$; neither $f$ itself nor $\gamma f$ is an
additional longitudinal mass flow.

Paired transfer gives $e+a=r$. In the invariant upstream free branch,
$S=C$, hence
\begin{equation}
 \dot e=-\frac vh\min\left\{e,
                     \frac{\kc(v)-r+e}{\gamma}\right\},\qquad a=r-e.
 \label{eq:h1time}
\end{equation}
The scalar right-hand side is Lipschitz. At $e=0$ it is zero, and the
receiving alternative vanishes at $a=\kc(v)$ if that boundary is reached.
Thus the initial bounds imply
\begin{equation}
 0\leq e(t)\leq e_0,\quad
 a_0\leq a(t)\leq\min\{r,\kc(v)\},\quad b_0\leq\kc(v).
 \label{eq:h1bounds}
\end{equation}
Under \eqref{eq:domain}, $r<R/v_+\leq R/v\leq\kc(v)$, with a strict
margin inherited from $v_+r<R$. Uniform source and common motion preserve
the remaining body's homogeneity provided the terminal entrance accepts it.
Section~\ref{sec:terminal} supplies that protection.

\subsection{Progress, explicit branches and the speed effect}
Put $z=vt$ and $\Delta(v)=\kc(v)-r>0$. Then
\begin{equation}
 e_z=-\frac1h\min\{e,(\Delta+e)/\gamma\}. \label{eq:h1progress}
\end{equation}
On a donor-limited interval starting at $(z_s,e_s)$,
\begin{equation}
 e(z)=e_s e^{-(z-z_s)/h}. \label{eq:h1donor}
\end{equation}
On a receiving-limited interval,
\begin{equation}
 e(z)=(e_s+\Delta)e^{-(z-z_s)/(\gamma h)}-\Delta.
 \label{eq:h1receiving}
\end{equation}
For $\gamma>1$ the switch is at $e=\Delta/(\gamma-1)$; the minimum
selects the branch. A receiving-limited solution decreases to that value
and then follows the donor branch. If $\gamma=1$, positive $\Delta$
makes the donor branch active throughout.

The explicit $v$ factor cancels in progress coordinates. Therefore travel
time alone no longer guarantees improved preparation. Since $\kc(v)$
decreases with $v$, the scalar comparison principle applied to
\eqref{eq:h1progress} makes $e(z;v)$ nondecreasing with advice at fixed
$z$, allowing equality on donor-only intervals. Speed dependence is
through the declared receiving headroom.

\section{Finite footprint, handoff arrivals and upstream mass balances}
\label{sec:handoff}
Early transfer stops at the fixed cutoff $L_p$. No ordinary lateral source
acts inside the terminal; residual E merges at $L$. With zero initial
terminal histories, entries satisfy almost everywhere
\begin{align}
 t_f^p&=(L_p-\ell)/v,\quad t_r^p=L_p/v,\nonumber\\
 \dot F_E&=ve(t),\quad \dot F_A=v(r-e(t)),\quad\dot F_B=vb_0
       \quad(t_f^p<t<t_r^p), \label{eq:entry}
\end{align}
and vanish outside that interval. Point values of flux at endpoints do
not affect continuous cumulative mass.

The remaining occupied preparation length is
\begin{equation}
 s(t)=\begin{cases}
 \ell,&0\leq t\leq t_f^p,\\
 L_p-vt,&t_f^p<t<t_r^p,\\
 0,&t\geq t_r^p.
 \end{cases}\qquad
 J_{\rm lat}(t)=\int_0^t s(u)f(u)/h\dd u. \label{eq:occupied}
\end{equation}
Differentiating interval by interval and matching initial values gives
\begin{align}
 s\rho_0+F_E+F_A+F_B&=M,\nonumber\\
 se+F_E+J_{\rm lat}&=\ell e_0,\nonumber\\
 sa+F_A-J_{\rm lat}&=\ell a_0,\qquad sb_0+F_B=\ell b_0.
 \label{eq:upstreammass}
\end{align}
An early-merged unit belongs to A once; it is not charged as a later forced
E merge. The uniform body's density state may be continued algebraically
after preparation empties, but no source or population then remains there.

\section{Terminal propagation, physical stock and storage protection}
\label{sec:terminal}
For each lane use minimum free travel $\tau=b/v$ and backward delay
$\tau_b=b/w$. Negative-time histories are zero. Define
\begin{equation}
 P_j(t)=F_j(t-\tau),\quad
 N_j(t)=F_j(t)-Y_j(t),\quad
 q_j(t)=P_j(t)-Y_j(t),\quad
 \lambda_j=\dot P_j,\quad g_j=\dot Y_j.
 \label{eq:terminalcounts}
\end{equation}
$N_j$ is physical stock, $q_j$ eligible unserved mass. Feasible departures
satisfy $0\leq Y_j\leq P_j\leq F_j$, so no unit exits before its minimum
travel. Queued demand is $C(v)$; an empty free approach offers $\lambda_j$.
The independent far lane has $Y_B=P_B$ under \eqref{eq:domain}.

For empty-link boundary receiving, the vacancy expression is
\begin{equation}
 V_j^{\rm vac}(t)=\begin{cases}
 k_jwt-F_j(t),&0<t\leq\tau_b,\\
 Y_j(t-\tau_b)+k_jb-F_j(t),&t>\tau_b.
 \end{cases} \label{eq:vacancy}
\end{equation}
By \eqref{eq:h1bounds}, $\dot F_j\leq C(v)<wk_j$. Initially
$F_j(t)\leq C(v)t<k_jwt$. Later
$V_j^{\rm vac}\geq k_jb-M>0$. Thus terminal entrance receiving never
holds preparation in the sufficient domain. This is finite physical
storage protection, not an infinite point buffer or a general spillback
solution. It also implies $N_j\leq M<k_jb$.

The eligible arrival pulse is
\begin{equation}
 t_f=(L-\ell)/v,\quad t_r=L/v,\qquad
 \lambda_A+\lambda_E=vr=:s_v<R,\quad
 \lambda'_A=-\lambda'_E\in[0,L_A],\quad L_A=v_+^2e_0/h
 \label{eq:pulse}
\end{equation}
inside the pulse. E arrivals decrease and A arrivals increase; only the
front and rear have external jumps. H1 branch switches preserve continuous
rates. This monotonic structure is essential to the later event proof.

\section{Current speed marks, immediate work and unique service roots}
\label{sec:root}
\subsection{Traffic-state closure and the fixed coefficient}
Free boundary marks use $u_j=v$; queued marks use $u_j=U(g_j)$. These
are equilibrium trace assumptions, not independent speed state or
whole-lane mean speeds. The selected W-F work expression is
\begin{equation}
 c_A=1,\qquad c_E=c_0+\chi g_A\pos{u_A-u_E},\qquad
 g_A+c_Eg_E\leq\budget. \label{eq:currentwork}
\end{equation}
It has no direct advice effect at fixed represented traffic and budget.
For queued E, with $a=g_A$, fixed A mark $u$ and available E work $x$,
\begin{equation}
 \He(g;a,u)=g\{c_0+\chi a\pos{u-U(g)}\}=x.
 \label{eq:root}
\end{equation}
On the active mismatch branch,
\begin{equation}
 \He_g=c_0+\chi a\{u-U(g)-gU'(g)\},\quad
 U'(g)=\frac{wK}{(K-g)^2},\quad
 gU'(g)\leq v(1+v/w)\quad(0\leq g\leq C(v)).
 \label{eq:rootderiv}
\end{equation}
Because $a\leq C(v)\leq C_\star$ and $v\leq v_\star$,
$\He_g\geq\eta>0$. Outside mismatch, $\He_g=c_0$; the continuous
function is strictly increasing across the kink. Also
$\He(0)=0$ and $\He(C(v))=c_0C(v)\geq\budget\geq x$.
There is exactly one feasible root. With the same bounds,
\begin{equation}
 c_0\leq c_E\leq c_{\max},\qquad
 x/c_{\max}\leq g_E\leq x/c_0,\qquad
 g_A+g_E\leq\budget\leq R\leq C(v).
 \label{eq:rootbounds}
\end{equation}
Thus sending/receiving ceilings and mass conservation are respected.

\subsection{What fixing the coefficient changed}
The earlier trial divided the mismatch term by $C(v)v$. Its multiplier
$\rho(v)=(w+v)/(Kv^2)$ has
$\rho'(v)=-(v+2w)/(Kv^3)<0$. When both approaches are queued at fixed
budget, A flow/mark are fixed, yet that trial changes E work with an
inactive advice cap. W-F removes that direct channel. It agrees with the
trial at the fixed anchor $v_+$; it does not preserve its advice derivative
away from that anchor.

On a regular free-A/queued-E branch with fixed $a,\budget$,
$u_A=v>U(g_E)$ and $x=\budget-a$. Implicit differentiation gives
\begin{equation}
 \frac{\partial g_E}{\partial v}
 =-\frac{\chi a g_E}{\He_g}\leq0. \label{eq:localbenefit}
\end{equation}
Lowering free A speed can therefore ease E entry in this local comparison.
This is not a full-cohort derivative: H1, arrivals, travel and recovery
also change. When both are queued below the advice cap, there is no direct
$v$ in the local root at fixed budget.

\section{Priority, cooperation and causal queue transitions}
\label{sec:allocation}
Free A claims its offer when it fits. Queued A uses equal-work sharing
with unused-work redistribution, including when E is free. A convenient
fixed-cost template is the minimal $z\geq0$ with
\begin{equation}
 \min(d_A,z)+\min(c_Ed_E,z)=\min(\budget,d_A+c_Ed_E),\quad
 g_A=\min(d_A,z),\quad g_E=\min(d_E,z/c_E).
 \label{eq:equalwork}
\end{equation}
The selected branches resolve the implicit cost as follows:
\begin{center}
\begin{tabular}{>{\raggedright\arraybackslash}p{27mm}>{\raggedright\arraybackslash}p{120mm}}
\toprule
Current modes & Service and feasibility\\
\midrule
A/E free & $g_A=\lambda_A$, $g_E=\lambda_E$ if
$\lambda_A+c_0\lambda_E\leq\budget$.\\[5pt]
A free, E queued & $g_A=a=\lambda_A\leq\budget$,
$\He(g_E;a,v)=\budget-a$.\\[5pt]
A/E queued & $g_A=\budget/2$,
$\He(g_E;\budget/2,U(\budget/2))=\budget/2$.\\[5pt]
A queued, E free & $g_E=\lambda_E$, $g_A=\budget-c_0\lambda_E$
when $c_0\lambda_E\leq\budget/2$; otherwise E has a strict offer
shortfall and enters its queued branch.\\
\bottomrule
\end{tabular}
\end{center}
Free E has $u_E=v$ and thus no positive mismatch term for either A mode.
With both queued, $g_E\leq g_A/c_0\leq g_A$; E's maneuver disadvantage
persists without a separate arbitrary flow-share parameter.

\textbf{Causal guards.} Initially both queues are zero and modes free.
An empty approach remains free while its offer is met. A strict free-mode
shortfall selects its queued trace; recomputed right-hand queue growth
must be compatible with nonnegative mass. While $q_j>0$,
\begin{equation}
 \dot q_j=\lambda_j-g_j,\qquad \dot Y_j=g_j. \label{eq:queuedynamics}
\end{equation}
At a true zero hit, retest free admission. Equality alone does not create
a queue; use maximal feasible free continuation and right-hand arrival
data at jumps. Retain recovery history across mode changes. No numerical
epsilon, branch latch, averaging rule or dwell time is added.

Why cooperation cannot be restricted literally to two positive queues:
at E clearance with A still queued, take unit costs, budget $4/5$,
$\lambda_A=3/5$, $\lambda_E=1/10$. Cooperative redistribution gives
$(g_A,g_E)=(7/10,1/10)$, so E stays clear. Strict all-state A priority
instead gives E no service at zero but $2/5$ at positive queue, requiring
$\dot q_E=1/10$ at zero and $-3/10$ above zero. No nonnegative absolutely
continuous continuation satisfies that literal toggle. The selected
current-state rule avoids it without allocation memory. This is a local
diagnostic, not a new primary initialization.

\section{A speed-marked recovery profile and recent-mass renewal}
\label{sec:recovery}
\subsection{The isolated boundary deduction}
Let $u_R=U(R)$ and $a_r>0$. Define
\begin{align}
 S(\theta;u)&=\sqrt{(w+u)^2+2wa_r\theta},\quad V=S-w,\nonumber\\
 Q_{\rm iso}(\theta;u)&=K\left[\theta-
                  \frac{S-w-u}{a_r}\right],\quad
 \partial_\theta Q_{\rm iso}=K(1-w/S)=c(V),\nonumber\\
 D(\theta;u)&=\pos{R-c(V(\theta;u))}. \label{eq:profile}
\end{align}
This reuses the isolated accelerated-insertion point expression to form
a capped boundary service loss. $V$ is not a vehicle's speed trajectory;
it parametrizes the boundary passing rate. The mark $u$ is the selected
representative insertion speed. The source isolated deduction does not
validate the aggregate overlap rule below.

For $u<u_R$,
\begin{align}
 \Theta(u)&=\frac{(w+u_R)^2-(w+u)^2}{2wa_r},\quad
 \Theta_{\max}=\Theta(0)=u_R/a_r+u_R^2/(2wa_r),\nonumber\\
 W(u)&=\int_0^\infty D(\theta;u)\dd\theta
 =\frac{K}{w+u_R}\frac{(u_R-u)^2}{2a_r}. \label{eq:profilearea}
\end{align}
For $u\geq u_R$, $D,W,\Theta$ are zero. To obtain the area, change
variable from age to $V$:
\begin{equation}
 \frac{d\theta}{dV}=\frac{w+V}{wa_r},\qquad
 D\frac{d\theta}{dV}=\frac{(K-R)(u_R-V)}{wa_r},
 \qquad \frac{K-R}{w}=\frac{K}{w+u_R}.
 \label{eq:areaderiv}
\end{equation}
Integrating from $u$ to $u_R$ gives \eqref{eq:profilearea}. On the active
part,
\begin{equation}
 D_\theta=-\frac{Kw^2a_r}{S^3},\quad
 D_u=-\frac{Kw(w+u)}{S^3},\quad
 D_{a_r}=-\frac{Kw^2\theta}{S^3}. \label{eq:profilederivs}
\end{equation}
Thus $D$ decreases with age and mark, $0\leq D\leq R$,
$|D_\theta|\leq a_rk_j$ and $|D_u|\leq k_j$. The bounds extend
through the continuous zero-loss cutoff. Physical age is not a numerical
step. $W$ is a lost-service area, not a second immediate E charge.

\subsection{The selected E-only renewal account}
For $Y=Y_E$, let $T(\xi)$ be its mass-rank admission-time inverse and
$a_m(t)=\pos{Y(t)-m}$. The selected account is
\begin{equation}
 \loss(t)=\frac1m\int_{a_m(t)}^{Y(t)}
 D(t-T(\xi);u_E(T(\xi)))\dd\xi,\qquad
 \budget(t)=\pos{R-\loss(t)}. \label{eq:renewal}
\end{equation}
Only E creates new loss. A influences E service/marks and consumes the
shared budget; A's slow passage is not independently charged as a new
profile. The account averages at most the latest mass $m$. New mass
retires equal oldest weight after fill while all profiles age. Recent
speed/age history, not just the reading, is retained predicted state.

Flat intervals of $Y$ carry no admission measure; inverse choices there
do not affect the result. Equivalently,
\begin{equation}
 \loss(t)=\frac1m\int_0^tD(t-s;u_E(s))g_E(s)
             \mathbf1_{\{Y(s)>a_m(t)\}}\dd s. \label{eq:renewaltime}
\end{equation}
Only ages up to $\Theta_{\max}$ and the latest mass $m$ can contribute.
Where derivatives exist, put $T_o=T(Y(t)-m)$ when $Y(t)>m$. Then
\begin{align}
 \dot\loss={}&\frac{g_E(t)}mD(0;u_E(t))
 +\frac1m\int_{a_m(t)}^{Y(t)}
 D_\theta(t-T(\xi);u_E(T(\xi)))\dd\xi\nonumber\\
 &-\frac{g_E(t)}m\mathbf1_{\{Y(t)>m\}}
 D(t-T_o;u_E(T_o)). \label{eq:renewalupdate}
\end{align}
These are injection, aging and oldest-weight retirement. There is no
division by an old flow rate. At fill and other endpoints use the a.e.
identity; the cumulative account is continuous. With $g_E=0$ only aging
remains. A fast new E entry can create zero loss and still retire an older
strong profile; this is a selected approximation requiring physical evidence.
